% Standalone problem document, generated from the adjacent README.md.
% Compile with XeLaTeX. Mathematical content is maintained in Markdown.
\documentclass[11pt,a4paper]{article}
\usepackage[fontsize=10.5pt]{scrextend}
\usepackage[margin=22mm,headheight=15pt,headsep=9mm,footskip=11mm]{geometry}
\usepackage{fontspec}
\setmainfont{texgyrepagella}[Extension=.otf,UprightFont=*-regular,BoldFont=*-bold,ItalicFont=*-italic,BoldItalicFont=*-bolditalic]
\setsansfont{texgyreheros}[Extension=.otf,UprightFont=*-regular,BoldFont=*-bold,ItalicFont=*-italic,BoldItalicFont=*-bolditalic]
\setmonofont{texgyrecursor}[Extension=.otf,UprightFont=*-regular,BoldFont=*-bold,ItalicFont=*-italic,BoldItalicFont=*-bolditalic,Scale=MatchLowercase]
\usepackage{amsmath,amssymb}
\usepackage{unicode-math}
\setmathfont{texgyrepagella-math.otf}
\usepackage{xcolor,microtype,enumitem,needspace,fancyhdr}
\usepackage{bookmark}
\definecolor{ink}{HTML}{17344B}
\definecolor{accent}{HTML}{197D80}
\definecolor{muted}{HTML}{596773}
\hypersetup{colorlinks=true,linkcolor=accent,urlcolor=accent,pdftitle={MI-03 - Sharp
additive triangle constant for an odd number of
contractions},pdfauthor={Open Problems in Numerical Linear Algebra}}
\urlstyle{same}
\setlength{\parindent}{0pt}
\setlength{\parskip}{5pt plus 1pt minus 1pt}
\linespread{1.04}
\setlength{\emergencystretch}{3em}
\widowpenalty=10000
\clubpenalty=10000
\displaywidowpenalty=10000
\allowdisplaybreaks[1]
\setlist{nosep,leftmargin=1.5em}
\providecommand{\tightlist}{\setlength{\itemsep}{0pt}\setlength{\parskip}{0pt}}
\providecommand{\pandocbounded}[1]{#1}
\pagestyle{fancy}
\fancyhf{}
\fancyhead[L]{\sffamily\footnotesize\color{muted}OPEN PROBLEMS IN NUMERICAL LINEAR ALGEBRA}
\fancyhead[R]{\sffamily\bfseries\footnotesize\color{accent}MI-03}
\fancyfoot[L]{\parbox[b]{0.9\textwidth}{\sffamily\fontsize{8}{10}\selectfont\color{muted}Editorial ratings; a bounded search cannot certify that a problem remains unsolved.\\Literature check: 2026-09-12}}
\fancyfoot[R]{\sffamily\footnotesize\color{muted}\thepage}
\renewcommand{\headrulewidth}{0.3pt}
\renewcommand{\headrule}{\hbox to\headwidth{\color{accent}\leaders\hrule height \headrulewidth\hfill}}
\makeatletter
\renewcommand{\section}{\Needspace{5\baselineskip}\@startsection{section}{1}{0pt}{12pt plus 2pt minus 2pt}{4pt}{\sffamily\bfseries\large\color{ink}}}
\renewcommand{\subsection}{\Needspace{4\baselineskip}\@startsection{subsection}{2}{0pt}{9pt plus 1pt minus 1pt}{3pt}{\sffamily\bfseries\normalsize\color{ink}}}
\makeatother
\setcounter{secnumdepth}{0}
\begin{document}
{\sffamily\small\color{muted}Matrix inequalities and norms\par}
\vspace{3pt}
{\raggedright\hyphenpenalty=10000\exhyphenpenalty=10000\sffamily\bfseries\fontsize{21}{24}\selectfont\color{ink}Sharp
additive triangle constant for an odd number of contractions\par}
\vspace{7pt}
{\sffamily\small\textbf{Difficulty:} challenging\quad\textbf{Importance:} interesting
to specialist\par}
{\sffamily\small\bfseries\color{accent}Status: Lean verified\par}
\vspace{4pt}
{\color{accent}\hrule height 0.8pt}
\vspace{6pt}
\textbf{Rating rationale:} Sharpness for every odd summand count needs a
new extremal construction or obstruction; the payoff is concentrated in
operator triangle inequalities.

\subsection{Resolution --- 2026-09-11}\label{resolution-2026-09-11}

\textbf{Affirmative result by Matthew J. Colbrook}, Department of
Applied Mathematics and Theoretical Physics, University of Cambridge.
\textbf{Independent proof review: PASS.}

The sharp additive contraction constant is \(c_k=k/4\) for every
\(k\ge2\), including every odd summand count. Exact dimension-two
extremizers match the universal upper bound.

The exact target is resolved. The original statement and source evidence
are retained below; its former difficulty rating is historical.

\textbf{Primary manuscript:}
\href{https://github.com/ajt60gaibb/OpenProblemsInNLA/blob/main/matrix-inequalities-and-norms/MI-03/solution.pdf}{complete
proof PDF},
\href{https://github.com/ajt60gaibb/OpenProblemsInNLA/blob/main/matrix-inequalities-and-norms/MI-03/solution.tex}{standalone
TeX}, Theorem 1.1 and its proof;
\href{https://github.com/ajt60gaibb/OpenProblemsInNLA/blob/main/matrix-inequalities-and-norms/MI-03/solution.md}{authorship
and scope}. The
\href{https://github.com/ajt60gaibb/OpenProblemsInNLA/blob/main/references/colbrook-matrix-2026-09-11/verification/reviews/MI-03-review.md}{independent
review} checks the full original argument and records its hash. The
draft was AI-assisted. That original review was informal; the later Lean
verification is documented below. External human peer review is not
claimed.
\href{https://github.com/ajt60gaibb/OpenProblemsInNLA/blob/main/references/colbrook-matrix-2026-09-11/README.md}{Submission
record}.

\subsection{Lean proof and verification evidence -
2026-09-12}\label{lean-proof-and-verification-evidence---2026-09-12}

\textbf{The complete original odd-summand conjecture is Lean verified.}
The
\href{https://github.com/sgstepaniants/OpenProblemsInNLA/tree/901ba5ffad3b57557b60c7360df67659d8b8aa21/matrix-inequalities-and-norms/MI-03/lean}{proof
at revision 901ba5f} proves that \(k/4\) is the least member of the full
admissible-constant set for every \(k\ge2\), then identifies its actual
real infimum. This includes every odd \(k\ge3\), with all positive
dimensions and every tuple of complex contractions. Genuine principal
CFC moduli and Euclidean operator norms are used; no literature
inequality is assumed as an unproved premise.

\textbf{Lean formalization:} George Stepaniants, Department of Computing
and Mathematical Sciences, California Institute of Technology, Pasadena,
California, USA, with AI-agent assistance. \textbf{Matthew J. Colbrook}
retains authorship of the mathematical proof and result;
\textbf{Jean-Christophe Bourin and Eun-Young Lee} retain the original
conjecture and prior-bound credit.

The eight
\href{https://github.com/sgstepaniants/OpenProblemsInNLA/blob/901ba5ffad3b57557b60c7360df67659d8b8aa21/matrix-inequalities-and-norms/MI-03/lean/Solution.lean}{checked
exports}, each with prefix \textit{NLA.MI03.}, are:

\begin{itemize}
\tightlist
\item
  \textit{modulus\_semantics}: genuine positive square root, square and
  norm identities.
\item
  \textit{contraction\_modulus}: positive modulus defects for every
  complex contraction.
\item
  \textit{positive\_decomposition}: the exact universal Gram-sum and
  positive-square identity.
\item
  \textit{universal\_upper\_bound}: admissibility of \(k/4\) for every
  \(k\ge2\).
\item
  \textit{root\_of\_unity\_data}: exact complex roots and their
  vanishing geometric sums.
\item
  \textit{sharpness\_witness}: true moduli, unit operator norms and full
  sums of the two-dimensional extremizers.
\item
  \textit{sharp\_constant}: actual least admissible constant and
  infimum, for every \(k\ge2\).
\item
  \textit{odd\_contraction\_conjecture}: the complete original
  assertion.
\end{itemize}

Two independent agents approved the
\href{https://github.com/ajt60gaibb/OpenProblemsInNLA/blob/main/matrix-inequalities-and-norms/MI-03/lean/reviews}{statements
and completed proof}.
\href{https://github.com/sgstepaniants/OpenProblemsInNLA/actions/runs/34722618003}{Linux
run 34722618003} matched all eight declarations with the sandboxed
Comparator and replayed the solution in Lean's default kernel. The
\href{https://github.com/ajt60gaibb/OpenProblemsInNLA/blob/main/matrix-inequalities-and-norms/MI-03/lean/verification/linux-2026-09-12}{original
artifacts and operational audit} bind all 173 input files and both
actual isolation/rejection-control suites. All 16
\href{https://github.com/ajt60gaibb/OpenProblemsInNLA/blob/main/matrix-inequalities-and-norms/MI-03/lean/verification/linux-2026-09-12/axiom-verification.json}{internal/public
transitive axiom reports} use exactly \textit{propext},
\textit{Classical.choice} and \textit{Quot.sound}. The operational
reviewer also served as final proof referee 1; these roles do not count
as a third mathematical referee. External human peer review is not
claimed.

The pins are \textbf{Lean 4.33.1},
\href{https://github.com/alerad/leancert/tree/621a43d7cf21f87872392a01e874f2f1dbddc926}{LeanCert
621a43d} and
\href{https://github.com/leanprover-community/mathlib4/tree/0df444a360eaa60ab8c11dca51a86af692955474}{Mathlib
0df444a}. LeanCert audits this exact proof's kernel trust; there is
\textbf{no numerical interval certificate}. Symbolic identities and
all-\(k\) roots avoid numerical phase approximation or interval
subdivision. The source's extra three-dimensional Hermitian extremizers
and rank classification are outside these exports; the complete original
target is covered. See the
\href{https://github.com/ajt60gaibb/OpenProblemsInNLA/blob/main/matrix-inequalities-and-norms/MI-03/lean/README.md}{project
guide},
\href{https://github.com/ajt60gaibb/OpenProblemsInNLA/blob/main/matrix-inequalities-and-norms/MI-03/lean/formalization.yaml}{manifest}
and
\href{https://github.com/ajt60gaibb/OpenProblemsInNLA/blob/main/matrix-inequalities-and-norms/MI-03/lean/lake-manifest.json}{dependency
pins}. From the immutable verified revision on a documented
\href{https://github.com/ajt60gaibb/OpenProblemsInNLA/blob/main/tools/lean/HARNESS.md}{non-root
Linux host}, reproduce with:

\begin{verbatim}
tools/lean/bootstrap.sh /absolute/path/to/nla-lean-tools
tools/lean/selftest.sh /absolute/path/to/nla-lean-tools
tools/lean/verify.sh \
  matrix-inequalities-and-norms/MI-03/lean \
  /absolute/path/to/nla-lean-tools
\end{verbatim}

\subsection{Problem statement}\label{problem-statement}

For each integer \(k\ge2\), let \(c_k\) be the infimum of all \(c\ge0\)
such that

\[\left|\sum_{j=1}^k A_j\right|\preceq cI_n+\sum_{j=1}^k|A_j|\]

for every \(n\ge1\) and all \(A_1,\ldots,A_k\in\mathbb C^{n\times n}\)
satisfying \(\|A_j\|_2\le1\). Here \(|A|=(A^*A)^{1/2}\), \(\|\cdot\|_2\)
denotes the operator norm, and \(X\preceq Y\) means \(Y-X\) is positive
semidefinite.

Is \(c_k=k/4\) for every odd integer \(k\ge3\)?

\subsection{Why it matters}\label{why-it-matters}

The scalar absolute-value triangle inequality fails in matrix order. The
optimal additive correction quantifies that failure for uniformly
bounded summands, in a form usable in positive block-matrix estimates.

\newpage
\subsection{References}\label{references}

\begin{enumerate}
\def\labelenumi{\arabic{enumi}.}
\tightlist
\item
  J.-C. Bourin and E.-Y. Lee, \emph{Diagonal and off-diagonal blocks of
  positive definite partitioned matrices}, arXiv:2307.02034v3 (15
  December 2023), Corollary 4.4 and the conjecture immediately following
  Remark 4.5. \href{https://arxiv.org/html/2307.02034}{Primary text}.
\item
  E.-Y. Lee, \emph{How to compare the absolute values of operator sums
  and the sums of absolute values?}, Operators and Matrices 6(3) (2012),
  613--619; §2 supplies related triangle inequalities.
  \href{https://files.ele-math.com/articles/oam-06-42.pdf}{Primary
  paper}, \href{https://doi.org/10.7153/oam-06-42}{DOI}.
\end{enumerate}

\subsection{Status check --- 2026-09-10}\label{status-check-2026-09-10}

Corollary 4.4 establishes \(c_k\le k/4\) for every \(k\) and sharpness
for even \(k\); the source separately conjectures sharpness for all odd
\(k>1\). Its latest version remains v3. Searches included
\textit{Bourin Lee contractions odd k constant k/4},
\textit{three contractions 3/4 sharp conjecture}, and
\textit{Bourin contractions sharp 2025 2026}. The 2026
symmetric-modulus papers concern different inequalities. No resolution
of this additive odd-summand problem was located.

\textbf{Audit update (2026-09-10):} Rechecked the conjecture after
Remark 4.5 in Bourin--Lee and searched for odd-summand sharpness
results. The proved even-summand statement is outside the target, which
already restricts to odd \(k\). This is a bounded literature check, not
a proof that no solution exists.
\end{document}
